\documentclass[11pt,letterpaper]{article}
\usepackage{../assets/scientific_report}
\usepackage{float}
\usepackage{listings}
\usepackage{xcolor}
\usepackage{graphicx}
\usepackage{geometry}
\geometry{headheight=23pt, top=2cm, bottom=2cm, left=2.5cm, right=2.5cm}
\setlength{\parskip}{0.5em}
\setlength{\parindent}{0pt}

\lstset{
    language=C,
    basicstyle=\ttfamily\small,
    keywordstyle=\color{blue}\bfseries,
    commentstyle=\color{green!60!black}\itshape,
    stringstyle=\color{red!80!black},
    numbers=left,
    numberstyle=\tiny\color{gray},
    numbersep=8pt,
    frame=single,
    breaklines=true,
    captionpos=b,
    tabsize=4,
    showstringspaces=false
}

\begin{document}
\pagestyle{plain}

\begin{center}
    {\LARGE\bfseries Processamento paralelo de lista encadeada com \texttt{task}}\\[0.4em]
    {\normalsize Eugênio Vitor Lopes dos Santos}\\[0.3em]
    {\small DCA3703 -- Programação Paralela, Universidade Federal do Rio Grande do Norte}\\[0.3em]
\end{center}

\section{Enunciado}
Implementar em C um programa que cria uma lista encadeada com nós contendo nomes de arquivos fictícios. Dentro de uma região paralela, o programa deve percorrer a lista e criar uma tarefa com \texttt{\#pragma omp task} para processar cada nó. Cada tarefa imprime o nome do arquivo e o identificador da thread que a executou (\texttt{omp\_get\_thread\_num}). Após executar o programa, deve-se refletir: todos os nós foram processados? Algum foi processado mais de uma vez ou ignorado? O comportamento muda entre execuções? Como garantir que cada nó seja processado uma única vez e por apenas uma tarefa?

\section{Fundamentação teórica}

\subsection{Modelo de tarefas do OpenMP}
O OpenMP oferece o modelo de tarefas (\texttt{task}) para expressar paralelismo em estruturas de dados irregulares, como listas encadeadas, árvores e grafos. Diferentemente de \texttt{parallel for}, que exige que o número de iterações seja conhecido antes da execução, as tarefas podem ser criadas dinamicamente conforme a estrutura é percorrida.

Uma diretiva \texttt{\#pragma omp task} delimita um bloco de código que pode ser executado por qualquer thread do time. A thread que criou a tarefa continua o fluxo normal; a tarefa entra em uma fila e qualquer thread pode executá-la. Se a thread criadora não tiver trabalho imediato e houver tarefas pendentes, ela pode executá-las. Caso contrário, threads ociosas podem ``roubar'' tarefas da fila de outras, mecanismo conhecido como \textit{work-stealing} (roubo de trabalho).

\subsection{As diretivas \texttt{single} e \texttt{task}}
\begin{itemize}
    \item \texttt{\#pragma omp single}: delimita um bloco que deve ser executado por exatamente uma thread do time. As demais threads saltam o bloco e prosseguem. É o mecanismo correto para que apenas uma thread percorra a lista e crie as tarefas.
    \item \texttt{\#pragma omp task}: cria uma tarefa com o bloco de código seguinte. A tarefa pode ser executada imediatamente ou postergada, dependendo do escalonador.
    \item \texttt{\#pragma omp taskwait}: barreira que aguarda a conclusão de todas as tarefas criadas dentro do escopo atual. É necessária quando o código após ela depende do resultado das tarefas.
\end{itemize}

\subsection{Escopo de variáveis em tarefas}
Uma variável referenciada dentro de uma \texttt{task} segue as mesmas regras de escopo da região \texttt{parallel} que a contém, a menos que uma cláusula de escopo seja explicitada:
\begin{itemize}
    \item \texttt{shared(var)}: todas as tarefas compartilham o mesmo endereço de memória. Se a variável muda após a criação da tarefa, a tarefa vê o novo valor, o que, no caso de um ponteiro de travessia, faz todas as tarefas enxergarem o mesmo nó (ou NULL).
    \item \texttt{firstprivate(var)}: cada tarefa recebe uma cópia privada inicializada com o valor da variável no momento da criação. É a cláusula adequada para capturar o ponteiro \texttt{atual} em cada iteração de um laço.
    \item \texttt{private(var)}: cria uma cópia privada sem inicialização (valor indefinido). Causa comportamento imprevisível.
\end{itemize}

\section{Implementação}

Os arquivos-fonte \texttt{e0.c} a \texttt{e5.c} apresentam o modelo de tarefas de forma incremental. As versões com erro e1 e e2 têm o mesmo problema; mantemos e2 por incluir o contador que quantifica as repetições. As fun\c c\~oes auxiliares (\texttt{criar\_no}, \texttt{inserir\_no\_fim}, \texttt{liberar\_lista}) e a estrutura \texttt{No} s\~ao id\^enticas em todas as vers\~oes; por isso, os trechos a seguir mostram apenas a \texttt{main()} de cada vers\~ao, onde as diretivas OpenMP e as cl\'ausulas de escopo variam. O c\'odigo completo de cada arquivo encontra-se no Ap\^endice~\ref{apendice:codigos}.

Compilamos os códigos-fonte com GCC 13.3, padrão C11, avisos ativados (\texttt{-Wall}) e suporte a OpenMP (\texttt{-fopenmp}).

\begin{lstlisting}[language=bash, caption={Compilação e execução dos programas}]
gcc -std=c11 -Wall -fopenmp e0.c -o e0
gcc -std=c11 -Wall -fopenmp e2.c -o e2
gcc -std=c11 -Wall -fopenmp e3.c -o e3
gcc -std=c11 -Wall -fopenmp e4.c -o e4
gcc -std=c11 -Wall -fopenmp e5.c -o e5

OMP_NUM_THREADS=4 ./e2
OMP_NUM_THREADS=4 ./e5
\end{lstlisting}

\subsection{Versão base sequencial (e0.c)}
Constrói a lista encadeada com 5 nomes de arquivo e a percorre sequencialmente, sem qualquer diretiva OpenMP. Serve como referência para o comportamento esperado.

\lstinputlisting[caption={e0.c, versão sequencial base}]{e0.c}

\subsection{Erro 1: ausência de \texttt{single} (e2.c)}
O laço de travessia fica diretamente dentro de \texttt{\#pragma omp parallel}. Todas as $T$ threads executam o laço completo e cada uma cria $N$ tarefas, totalizando $N \times T$ tarefas e $N \times T$ processamentos do mesmo nó. O contador \texttt{processados} quantifica o excesso.

\begin{lstlisting}[caption={main() de e2.c, parallel sem single}]
int main() {
    No *cabeca = NULL;
    /* cria lista com 5 n00f3s */

        int nos_lista = 5;
    int processados = 0;

    #pragma omp parallel
    {
        No *atual = cabeca;
        while (atual != NULL) {
            #pragma omp task
            {
                printf("  [Thread %d] %s\n",
                       omp_get_thread_num(), atual->nome_arquivo);
                processados++;
            }
                atual = atual->next;
            }
        }
        #pragma omp taskwait
    }

    printf("Nos na lista (N): %d\n"
           "Processamentos totais (Y): %d\n"
           "Relacao: Y = N x T = %d x %d = %d (esperado)\n",
           nos_lista, processados,
           nos_lista, (processados / nos_lista), processados);

    liberar_lista(cabeca);
    return 0;
}
\end{lstlisting}

\subsection{Erro 2: \texttt{single} sem \texttt{firstprivate} (e3.c)}
Há \texttt{\#pragma omp single} para que apenas uma thread crie as tarefas. A cláusula \texttt{shared(atual)} faz todas as tarefas compartilharem o mesmo ponteiro. Como o laço avança o ponteiro imediatamente após criar cada tarefa, as tarefas executadas posteriormente leem um valor já modificado (outro nó ou NULL).

\begin{lstlisting}[caption={main() de e3.c, single + shared}]
int main() {
    No *cabeca = NULL;
    inserir_no_fim(&cabeca, "doc1.txt");

    int nos_lista = 1;
    int processados = 0;

    #pragma omp parallel
    {
        #pragma omp single
        {
            No *atual = cabeca;
            while (atual != NULL) {
                #pragma omp task shared(atual)
                {
                    printf("  [Thread %d] atual aponta para: %s\n",
                           omp_get_thread_num(),
                           atual ? atual->nome_arquivo : "(NULL)");
                    processados++;
                }
                atual = atual->next;
            }
        }
    }

    printf("Nos na lista: %d | Processamentos contados: %d\n",
           nos_lista, processados);

    liberar_lista(cabeca);
    return 0;
}
\end{lstlisting}

\subsection{Erro 3: \texttt{firstprivate} sem \texttt{taskwait} (e4.c)}
As tarefas recebem \texttt{firstprivate(atual)} e cada uma captura o endereço correto do nó — o nome impresso é sempre correto. No entanto, sem \texttt{taskwait}, não há garantia formal de que as tarefas estejam concluídas antes que código subsequente dependa de seus resultados. A barreira implícita no fim da região \texttt{parallel} compensa neste caso específico, mas não é a forma portável de sincronizar.

\begin{lstlisting}[caption={main() de e4.c, firstprivate sem taskwait}]
int main() {
    No *cabeca = NULL;
    inserir_no_fim(&cabeca, "doc1.txt");

    int nos_lista = 1;
    int processados = 0;

    #pragma omp parallel
    {
        #pragma omp single
        {
            No *atual = cabeca;
            while (atual != NULL) {
                #pragma omp task firstprivate(atual)
                {
                    printf("  [Thread %d] Arquivo: %s\n",
                           omp_get_thread_num(), atual->nome_arquivo);
                    processados++;
                }
                atual = atual->next;
            }
        }
    }

    printf("Nos na lista: %d | Processamentos contados: %d\n",
           nos_lista, processados);

    liberar_lista(cabeca);
    return 0;
}
\end{lstlisting}

\subsection{Versão correta (e5.c)}
A solução combina três elementos:
\begin{enumerate}
    \item \texttt{single}: apenas uma thread percorre a lista e cria as tarefas;
    \item \texttt{firstprivate(atual)}: cada tarefa captura o endereço do nó no momento da criação;
    \item \texttt{taskwait}: sincroniza o término das tarefas antes de prosseguir.
\end{enumerate}

\begin{lstlisting}[caption={main() de e5.c, versão correta}]
int main() {
    No *cabeca = NULL;
    inserir_no_fim(&cabeca, "doc1.txt");

    int nos_lista = 1;
    int processados = 0;

    #pragma omp parallel
    {
        #pragma omp single
        {
            No *atual = cabeca;
            while (atual != NULL) {
                #pragma omp task firstprivate(atual)
                {
                    printf("  [Thread %d/%d] Arquivo: %s\n",
                           omp_get_thread_num(),
                           omp_get_num_threads(),
                           atual->nome_arquivo);
                    processados++;
                }
                atual = atual->next;
            }
        }
    }

    printf("Nos na lista: %d | Processamentos: %d (devem ser iguais!)\n",
           nos_lista, processados);

    liberar_lista(cabeca);
    return 0;
}
\end{lstlisting}

\newpage
\section{Resultados e discussão}

A Tabela~\ref{tab:comparacao} resume o comportamento de cada versão com 4 threads.

\begin{table}[H]
\centering
\caption{Comparação do processamento das versões com 4 threads}
\label{tab:comparacao}
\begin{tabular}{llll}
\hline
\textbf{Arquivo} & \textbf{Nós process.} & \textbf{Nome correto?} & \textbf{Problema} \\ \hline
\texttt{e0.c} & 5 (sequencial) & Sim & Sem paralelismo \\
\texttt{e2.c} & 20 ($N\times T$) & Sim & Contador quantifica o excesso \\
\texttt{e3.c} & 5 & Sim (por acaso) & \texttt{shared(current)}: ponteiro inválido \\
\texttt{e4.c} & 5 & Sim & Sem \texttt{taskwait}: sem sincronização formal \\
\texttt{e5.c} & 5 & Sim & Correto: \texttt{single+firstprivate+taskwait} \\ \hline
\end{tabular}
\end{table}

\subsection{Ausência de \texttt{single} (e2.c)}
Com 4 threads e 5 nós, \texttt{e2.c} confirma \texttt{Y = 20}:

\begin{lstlisting}[language=bash, caption={Saída de e2.c com 4 threads confirmando $N \times T$}]
$ OMP_NUM_THREADS=4 ./e2
[Thread 0] doc1.txt
[Thread 0] img2.png
[Thread 0] dat3.csv
[Thread 0] rep4.pdf
[Thread 0] bkp5.zip
[Thread 0] doc1.txt
[Thread 0] img2.png
[Thread 0] dat3.csv
[Thread 3] bkp5.zip
[Thread 3] doc1.txt
[Thread 3] img2.png
[Thread 3] dat3.csv
[Thread 3] rep4.pdf
[Thread 3] bkp5.zip
[Thread 0] rep4.pdf
[Thread 3] doc1.txt
[Thread 0] img2.png
[Thread 3] dat3.csv
[Thread 1] rep4.pdf
[Thread 1] bkp5.zip
Nos na lista (N): 5
Processamentos totais (Y): 20
Relacao: Y = N x T = 5 x 4 = 20 (esperado)
\end{lstlisting}

Cada thread executou o laço independentemente e gerou 5 tarefas. O programa gastou $T$ vezes mais trabalho que o necessário. Se as tarefas tivessem efeitos colaterais (escrita em arquivo, atualização de banco de dados), cada nó seria alterado $T$ vezes.

\subsection{Compartilhamento do ponteiro (e3.c)}
Mesmo com \texttt{single}, se a cláusula de escopo é \texttt{shared(current)}, o ponteiro avança enquanto as tarefas ainda não executaram. O comportamento é não-determinístico: o contador acusa o número correto de nós, mas algumas tarefas podem imprimir NULL ou nomes trocados, dependendo do instante em que o escalonador as ativa.

Sem \texttt{firstprivate}, a correção do resultado depende de uma condição de corrida favorável. Não há garantia alguma. Comparando \texttt{e3.c} (shared) com \texttt{e4.c} (firstprivate), vê-se que a cláusula de escopo é decisiva.

\subsection{Ausência de \texttt{taskwait} (e4.c)}
Com \texttt{firstprivate}, \texttt{e4.c} imprime os 5 nomes corretos — cada tarefa captura o valor do ponteiro \texttt{atual} no instante da criação, então o laço \texttt{atual = atual->next} não afeta tarefas já criadas. No entanto, sem \texttt{taskwait}, não há sincronização formal entre a criação das tarefas e código subsequente que dependa delas. Neste caso, a barreira implícita no fim da região \texttt{parallel} compensa, mas o uso de \texttt{taskwait} é a forma correta e portável de garantir a conclusão.

\subsection{Versão correta}

Com \texttt{single}, \texttt{firstprivate(atual)} e \texttt{taskwait} (código análogo ao de \texttt{e4.c} acrescentando a barreira), cada nó é processado exatamente uma vez, como mostra a saída de \texttt{e5.c}:

\begin{lstlisting}[language=bash, caption={Execução da versão correta com 4 threads}]
$ OMP_NUM_THREADS=4 ./e5
[Thread 1/4] Arquivo: doc1.txt
[Thread 1/4] Arquivo: bkp5.zip
[Thread 2/4] Arquivo: img2.png
[Thread 3/4] Arquivo: dat3.csv
[Thread 0/4] Arquivo: rep4.pdf
Nos na lista: 5 | Processamentos: 5 (devem ser iguais!)
\end{lstlisting}

O programa imprimiu 5 linhas (uma por nó). Quatro threads diferentes executaram as tarefas. A ordem varia entre execuções, mas a contagem permanece 5.

\section{Conclusão}

O paralelismo de listas encadeadas com OpenMP exige três elementos:

\begin{enumerate}
    \item \texttt{single}: restringe a criação de tarefas a uma única thread, evitando a multiplicação de trabalho;
    \item \texttt{firstprivate}: captura o valor do ponteiro de travessia no momento da criação, impedindo que tarefas compartilhem um endereço que avança;
    \item \texttt{taskwait}: sincroniza o término das tarefas antes de prosseguir.
\end{enumerate}

A ausência de qualquer um desses elementos resulta em um dos três erros: (i) cada nó é processado $T$ vezes (sem \texttt{single}), (ii) tarefas leem ponteiros inválidos ou trocados (sem \texttt{firstprivate}) ou (iii) não há garantia de conclusão antes de código subsequente (sem \texttt{taskwait}). O comportamento varia entre execuções porque o escalonamento de tarefas é não-determinístico: a thread que executa cada tarefa, a ordem de execução e o instante de ativação dependem do estado do runtime.

\newpage
\appendix
\section{Código-fonte completo}
\label{apendice:codigos}

\lstinputlisting[caption={e0.c, versão sequencial base}]{e0.c}

\lstinputlisting[caption={e2.c, parallel sem single com contador}]{e2.c}

\lstinputlisting[caption={e3.c, single sem firstprivate}]{e3.c}

\lstinputlisting[caption={e4.c, firstprivate sem taskwait}]{e4.c}

\lstinputlisting[caption={e5.c, versão correta com single + firstprivate + taskwait}]{e5.c}

\end{document}